\documentclass{article}
\usepackage[utf8]{inputenc}
\usepackage[T1]{fontenc}
\usepackage{geometry}
\geometry{
    a4paper,
    left=2.5cm,
    right=2.5cm,
    top=2.5cm,
    bottom=2.5cm
}
\usepackage{graphicx}
\usepackage{float}
\usepackage{hyperref}
\usepackage[english]{babel}
\usepackage{amsmath}
\usepackage{amssymb}
\usepackage{array}
\usepackage{booktabs}
\usepackage{caption}
\usepackage{subcaption}
\usepackage{enumitem}
\usepackage{fancyhdr}
\usepackage{titlesec}
\usepackage{url}
\usepackage{textcomp}
\usepackage{xcolor}
\usepackage{listings}
\usepackage{multirow}
\usepackage{longtable}
\usepackage{tabularx}
\usepackage{times}

% Set up page style
\pagestyle{fancy}
\fancyhf{}
\rfoot{\thepage}
\renewcommand{\headrulewidth}{0pt}

% Title formatting
\titleformat{\chapter}[display]{\normalfont\bfseries}{\chaptertitle\thechapter}{0pt}{\Huge}
\titlespacing*{\chapter}{0pt}{50pt}{40pt}

% Document info
\title{Template Document}
\author{Academic Support Agent}
\date{\today}

\begin{document}
\maketitle

% D but du document
\thispagestyle{empty}	                        % Pour  viter d'avoir un en-t te et un pied de page sur la page couverture
\begin{center}
    \includegraphics[width=5cm]{logo.png}	    % Pour inclure le logo (on pr cise la largeur de l'image)
\end{center}
\vspace{4cm}	                                % Espacement vertical
\begin{center}	                                % On centre le texte
    {\huge \bf Titre }\\                        % \huge fait que le texte est gros, \bf fait que le texte est gras
    \vspace{3cm}
    \large Travail pr sent    M. Untel dans le cadre du cours \emph{XXX-YYYY : Nom du cours} \\
    \vspace{3cm}
    Travail r alis  par A et B\\
    \vfill	                                    % On va jusqu'au bas de la page avant de mettre le texte ci-dessous
    date de remise
\end{center}
\newpage
\begin{question}
    Ceci est la premi re question. Ceci est la premi re question. Ceci est la premi re question. Ceci est la premi re question. Ceci est la premi re question.\\
    
    Pour y r pondre, je dois  crire des  quations :
    \begin{align*}
        a&=b\\
        c&=\sqrt{x}
    \end{align*}
\end{question}
\begin{question}
    Dans cette question, nous devons trouver une photo du p re des processus stochastiques. La voici:
    \begin{figure}[!h]      % Le !h entre crochet c'est pour forcer la position de la photo ici.
        \centering
        \includegraphics[width=10cm]{AAMarkov.jpg}
        \caption{Andre  Andre evitch Markov (1856-1922) (voir \cite{wiki:Markov})}
        \label{fig:markov}
    \end{figure}
\end{question}
Voici un tableau.
\begin{table}[h!]
    \centering
    \begin{tabular}{|c|c|c|c|}
        \hline
          & 1    & 2    & 3 \\ \hline
        1 & 1    & 2    & 3 \\ \hline
        2 & 2    & 4    & 6 \\ \hline
    \end{tabular}
\end{table}

\newpage
    Voici une matrice identit :
    \[
        \begin{pmatrix}
         1    &0    & 0 \\
         0    & 1    & 0 \\
         0    & 0    & 1 \\
        \end{pmatrix}
    \]
    \begin{center}
        Voici du texte au centre.
    \end{center}
    \begin{flushright}
                {\color{red}Voici du texte   droite en rouge.}
    \end{flushright}
Pour en apprendre sur la th orie des jeux, vous pouvez consulter l'ouvrage de James Webb (voir \cite{Webb}).
\newpage
\end{document}